\documentclass[10pt,a4paper]{amsart}
\usepackage[margin=19mm]{geometry}
\usepackage{amsmath,amssymb,mathtools}
\usepackage[scr=rsfs,cal=euler]{mathalfa}
\usepackage{xfrac}
\usepackage[unicode,hidelinks,bookmarks=false]{hyperref}
\newcommand{\ie}{i.e.\ }
\newcommand{\cf}{cf.\ }
\newcommand{\resp}{respectively\ }
\newcommand{\Subsection}{\subsection}
\providecommand{\nobreakdash}{\nobreak\hbox{-}}
\setlength{\emergencystretch}{2em}
\multlinegap0pt
\usepackage{amssymb}
\newtheorem{corollary}{Corollary}
\newtheorem{proposition}{Proposition}
\title{Discrete Gaussian free fields on Hamming graphs}
\author{Shuhei Mano}
\date{Source: arXiv:2512.24199v2 (CC BY 4.0)}
\begin{document}
\maketitle

\noindent\emph{Source and licence.} Discrete Gaussian free fields on Hamming graphs. Version arXiv:2512.24199v2 (CC BY 4.0), distributed by arXiv under the Creative Commons Attribution 4.0 International licence, http://creativecommons.org/licenses/by/4.0/, as recorded in the arXiv licence metadata for this exact version and shown on the abstract page https://arxiv.org/abs/2512.24199v2. Full licence notice: \texttt{licenses/arxiv-2512.24199-CC-BY-4.0.txt}.\par\smallskip

\noindent\textit{Editorial scope.} Counted unit: the article's proposition pro:limcovd (source0.tex lines 1478--1488). The article's complete original proof of it is delivered with it (source0.tex lines 1490--1628, one \texttt{proof} environment of 139 lines). No mathematics is written, repaired or re-derived: the statement and the proof are the article's own lines, in the article's own notation, and no retained line is substituted or re-flowed.  Retained uncounted as the source-local context the proof needs: lines 161--166; lines 606--608; lines 621--626; lines 672--679; lines 1281--1309.  Nothing was omitted from any retained span, so the package carries no omission record.  No retained line contains a citation, so the wrapper carries no bibliography and no bibliography entry.  No retained line contains a figure environment, an image-inclusion command or drawing code, so no image is delivered and the package declares no asset dependency.  Editorial formatting only, in addition: an AMS \texttt{amsart} preamble that restates the article's own declarations and macros used by the retained lines, namely the environments \texttt{corollary}, \texttt{proposition} on separate counters, together with the standard packages those lines need.  The retained environments print in this document as Corollary 1, Proposition 1 in order of appearance, whereas the article prints the same items with the numbering of its own full document.  The complete native source of the article as distributed in the arXiv e-print for this exact version is bundled unchanged as \texttt{sources/191861/source0.tex}.
\begin{equation}\label{gibbs}
  \mu_\mathcal{Y}(S)=\frac{1}{Z_\mathcal{Y}}\int_S
  e^{-\beta H_\mathcal{Y}\{(\xi_x)_{x\in\mathcal{X}}\}}
  \prod_{x\in\mathcal{Y}} d\xi_x
  \quad \text{for~each} \quad S\in\mathcal{B},
\end{equation}

\begin{equation}\label{gen}
  \sum_{l=0}^dK_l(j)s^l=(1+(n-1)s)^{d-j}(1-s)^j.
\end{equation}

\begin{equation}\label{Kbound}
  \frac{|K_i(j)|}{\kappa_i}\le
  \sum_{l=\max(0,i+j-1)}^{\min(i,j)}(n-1)^{-l}
  \frac{\binom{j}{l}\binom{d-j}{i-l}}{\binom{d}{i}}
  =\mathbf{E}[(n-1)^{-{L_{i,j}}}]\le 1,
\end{equation}

\begin{align}
  G_{\alpha,\emptyset}(x,x')
  &=\sum_{t=0}^\infty \alpha^t P_t(\rho(x,x'))
  =\frac{1}{n^d}\sum_{t=0}^\infty\sum_{j=0}^d
  (\alpha\lambda_j)^t K_j(\rho(x,x'))\nonumber\\
  &=\frac{1}{n^d}\sum_{j=0}^d\frac{K_j(\rho(x,x'))}
  {1-\alpha\lambda_j}. \label{green_h}
\end{align}

\begin{corollary}\label{cor:lim}
  Set $m>0$ and $\partial=\emptyset$. The free energy of
  a Gaussian free field given by the Gibbs measure
  \eqref{gibbs} per particle satisfies the following.
  \begin{itemize}
  \item[$(i)$]     
  Assume there exists $\lambda_*=\lim_{d\to\infty}\lambda_{\lfloor(1-1/n)d\rfloor}$, and that, for every $\epsilon>0$, there exists
  $\delta>0$ such that
  $\max_{i\in A_d(\epsilon)}|\lambda_i-\lambda_*|\to 0$ as
  $d\to\infty$, where
  $A_d(\epsilon)=\{i\in\{0,1,\ldots,d\}:|i-(1-1/n)d|<\epsilon d^{1/2+\delta}\}$.
  Then,
  \begin{equation*}
    -\frac{1}{\beta}
    \lim_{d\to\infty}\frac{\log Z_{\mathcal{X}}}{n^d}
    =\frac{1}{2\beta}\left\{\log(1+m^2-\lambda_{*})
    -\log\left(\frac{2\pi}{\beta}\right)\right\};
  \end{equation*}
  \item[$(ii)$]
  Assume there exists
  $\lambda_{\#}=\lim_{n\to\infty}\lambda_d$. Then,
  \begin{equation*}
    -\frac{1}{\beta}
    \lim_{n\to\infty}\frac{\log Z_{\mathcal{X}}}{n^d}
    =\frac{1}{2\beta}\left\{\log(1+m^2-\lambda_\#)-
    \log\left(\frac{2\pi}{\beta}\right)\right\}.
  \end{equation*}
  \end{itemize}
\end{corollary}  

\begin{proposition}\label{pro:limcovd}
  Let $\lambda_*$ be the constant defined in Corollary~\ref{cor:lim}.
  Set $m>0$ and $\partial=\emptyset$.
  For $(x,x')\in\mathcal{X}\times\mathcal{X}$, the covariance
  of a Gaussian free field $(g_x)_{x\in\mathcal{X}}$ given
  by the Gibbs measure \eqref{gibbs} satisfies
  \begin{equation*}
    \lim_{d\to\infty}{\rm Cov}(g_x,g_{x'})=\frac{1}{\beta}
    \frac{\delta_{x,x'}}{1+m^2-\lambda_*}.
  \end{equation*}
\end{proposition}

\begin{proof}
  If $x=x'$, the Green function \eqref{green_h} gives
  \begin{align*}
    {\rm Var}(g_x)
    &=\frac{\alpha}{\beta}G_{\alpha,\emptyset}(x,x)
    =\frac{1}{\beta n^d}\sum_{i=0}^d\frac{K_i(0)}
    {1/\alpha-\lambda_i}
    =\frac{1}{\beta n^d}\sum_{i=0}^d\frac{\kappa_i}
    {1/\alpha-\lambda_i}\\
    &=\frac{1}{\beta}\sum_{i=0}^d
    {d \choose i}\left(1-\frac{1}{n}\right)^i
    \left(\frac{1}{n}\right)^{d-i}
    \frac{1}{1/\alpha-\lambda_i}.
  \end{align*}
  The same argument used in the proof of
  Corollary~\ref{cor:lim} gives
  \[
  \lim_{d\to\infty}{\rm Var}(g_x)
  =\frac{1}{\beta}\frac{1}{1+m^2-\lambda_*}.
  \]
  If $x\neq x'$, we have
  \begin{equation*}
    \alpha G_{\alpha,\emptyset}(x,x')=\frac{1}{n^d}
    \sum_{i=0}^d \frac{K_i(\rho)}{1/\alpha-\lambda_i}
    =\sum_{i=0}^d
    \binom{d}{i}\left(1-\frac{1}{n}\right)^i
    \left(\frac{1}{n}\right)^{d-i}
    \frac{K_i(\rho)/\kappa_i}{1+m^2-\lambda_i},
  \end{equation*}
  where $\rho=\rho(x,x')$.
  Substituting $s=1$ into \eqref{gen}, we have $\sum_{i=0}^d K_i(\rho)=0$
  for $\rho>0$. Therefore, we may write
  \begin{align*}
    \alpha G_{\alpha,\emptyset}(x,x')
    &=\sum_{i=0}^d
    \binom{d}{i}\left(1-\frac{1}{n}\right)^i
    \left(\frac{1}{n}\right)^{d-i}
    \frac{K_i(\rho)}{\kappa_i}
    \left(\frac{1}{1+m^2-\lambda_i}-\frac{1}{1+m^2-\lambda_*}\right)\\
    &=\mathbf{E}\left[
    \frac{K_{I_d}(\rho)}{\kappa_{I_d}}
    \left(\frac{1}{1+m^2-\lambda_{I_d}}-\frac{1}{1+m^2-\lambda_*}
    \right)\right],
  \end{align*}
  where $I_d$ follows a binomial distribution. Then,
  \begin{align*}
    |\alpha G_{\alpha,\emptyset}(x,x')|&\le
    \mathbf{E}\left[\left|
    \frac{1}{1+m^2-\lambda_{I_d}}-\frac{1}{1+m^2-\lambda_*}
    \right|\right]\\
    &\le
    \max_{i\in A_d(\epsilon)}
    \left|\frac{1}{1+m^2-\lambda_{I_d}}-\frac{1}{1+m^2-\lambda_*}\right|
    +C'\mathbf{P}(I_d\in A_d^c(\epsilon))
  \end{align*}
  for a positive constant $C'$ independent of $d$, where we used
  \eqref{Kbound} in the first inequality, and the bound
  $-1/(n-1)\le \lambda_i \le 1$, $i\in\{0,1,\ldots,d\}$ in
  the second inequality. By using
  $\max_{i\in A_d(\epsilon)}|1/(1+m^2-\lambda_{I_d})-1/(1+m^2-\lambda_*)|\to 0$
  as $d\to\infty$ and
  $\mathbf{P}(I_d\in A_d^c(\epsilon))\le e^{-2\epsilon^2 d^{2\delta}}$
  by Hoeffding's inequality, we have
  $\lim_{d\to\infty}|\alpha G_{\alpha,\emptyset}(x,x')|=0$,
  where $A_d(\epsilon)$ is defined in Corollary~\ref{cor:lim}.
  Therefore, we conclude that
  $\lim_{d\to\infty}{\rm Cov}(g_x,g_{x'})=\lim_{d\to\infty}\alpha G_{\alpha,\emptyset}(x,x')/\beta=0$ for $x\neq x'$.
  %Since $-1/(n-1)\le \lambda_i\le 1$, $i\in\{0,\ldots,d\}$,
  %we have
  %\begin{equation}\label{pro:limcovd1}
  %  |\alpha G_{\alpha,\emptyset}(x,x')|\le
  %  \sum_{i=0}^d
  %  \binom{d}{i}\left(1-\frac{1}{n}\right)^i
  %  \left(\frac{1}{n}\right)^{d-i}
  %  \frac{|K_i(\rho)/\kappa_i|}{1+m^2-\lambda_i}
  %  \le\frac{1}{m^2}
  %  \sum_{i=0}^d
  %  \binom{d}{i}\left(1-\frac{1}{n}\right)^i
  %  \left(\frac{1}{n}\right)^{d-i}
  %  \frac{|K_i(\rho)|}{\kappa_i}.
  %\end{equation}
  %The proof is split into two cases:
  %fixed $\rho$ and $\rho=\rho_d\to\infty$ as $d\to\infty$.\\
  %\color{red}
  %(Fixed $\rho$) Assume $\rho$ is fixed. For $i\asymp d$, set
  %$i=ud$, $u\in(0,1]$. We have
  %\begin{align*}
  %  \frac{K_{ud}(\rho)}{\kappa_{ud}}&=
  %       {}_2F_1\{-ud,-\rho;-d;n/(n-1)\}
  %       =\sum_{l=0}^{\rho}
  %       \frac{(-ud)_l(-\rho)_l}
  %       {(-d)_l}\frac{(n/(n-1))^l}{l!}
  %       =\sum_{l=0}^{\rho}
  %       \left(\frac{un}{n-1}\right)^l
  %       \frac{(-\rho)_l}{l!}\{1+O(d^{-1})\}\\
  %       &=\sum_{l=0}^\rho\binom{\rho}{l}\left(-\frac{un}{n-1}\right)^l
  %       +O(d^{-1})=
  %       \left(1-\frac{un}{n-1}\right)^\rho+O(d^{-1}),
  %\end{align*}
  %where the index of the Krawtchouk polynomial is extended to
  %$(0,d]$ via the Gauss hypergeometric polynomial ${}_2F_1(a,b;c;z)$,
  %and we set $(a)_l=a(a+1)\cdots (a+l-1)$. Then, we have
  %\begin{equation}\label{pro:limcovd2}
  %  \sup_{i\in B_d}\left|\frac{|K_i(\rho)|}{\kappa_i}-\frac{|K_{(1-1/n)d}(\rho)|}{\kappa_{(1-1/n)d}}\right|
  %  =O(d^{-1/3}) \to 0, \quad d\to\infty,
  %\end{equation}
  %for $B_d:=\{i\in\{0,\ldots,d\}:|i-(1-1/n)d|\le\epsilon d^{2/3}$.
  %The right-hand side of \eqref{pro:limcovd1} is
  %$\mathbf{E}[|K_{I_d}(\rho)|/\kappa_{I_d}]/m^2$,
  %where $I_d$ follows the binomial distribution. Then,
  %\[
  %\mathbf{E}\left[\left|\frac{|K_{I_d}(\rho)|}{\kappa_{I_d}}-
  %\frac{|K_{(1-1/n)d}(\rho)|}{\kappa_{(1-1/n)d}}\right|\right]
  %\le
  %\sup_{i\in B_d}\left[\left|\frac{|K_{I_d}(\rho)|}{\kappa_{I_d}}-
  %\frac{|K_{(1-1/n)d}(\rho)|}{\kappa_{(1-1/n)d}}\right|\right]+
  %\mathbf{P}(I_d\in B_d^c),
  %\]
  %where we used \eqref{Kbound}. By using \eqref{pro:limcovd2} and
  %$\mathbf{P}(I_d\in B_d^c)\le e^{-2\epsilon^2d^{1/3}}$ which follows
  %by Hoeffding's inequality, we have
  %\[
  %\lim_{d\to\infty}\mathbf{E}\left[\frac{|K_{I_d}(\rho)|}{\kappa_{I_d}}\right]
  %=\lim_{d\to\infty}\frac{|K_{(1-1/n)d}(\rho)|}{\kappa_{(1-1/n)d}}=0.
  %\]
  %Therefore, we conclude that
  %$\lim_{d\to\infty}{\rm Cov}(g_x,g_{x'})=\lim_{d\to\infty}\alpha G_{\alpha,\emptyset}(x,x')/\beta=0$.\\
  %($\rho_d\to\infty$) The right hand-side of \eqref{pro:limcovd1} is
  %\[
  %\frac{1}{m^2n^d}
  %\sum_{i=0}^d\sum_{l=\max(0,i+\rho_d-d)}^{\min(i,\rho_d)}
  %(n-1)^{i-l}\binom{\rho_d}{l}\binom{d-\rho_d}{i-l}
  %=\frac{1}{m^2n^d}
  %\sum_{l=0}^\rho\binom{\rho_d}{l}\sum_{i=0}^{d-\rho_d}
  %\binom{d-\rho_d}{i}(n-1)^{i}=\frac{1}{m^2}\left(\frac{2}{n}\right)^{\rho_d}.
  %\]
  %Therefore, if $n\ge 3$, $\alpha G_{\alpha,\emptyset}(x,x')\to 0$
  %as $\rho_d\to\infty$ with $d\to\infty$.  
\end{proof}

\end{document}
